\section{Contoh Program}
\label{sec:contoh-program}

Dua contoh berikut memperlihatkan bagaimana nilai segmen dan offset
diperoleh dari dalam program, lalu bagaimana keduanya dipakai untuk
menghitung alamat fisik. Keduanya memakai prosedur bantu
\code{cetak\_hex} yang mencetak isi register \reg{AX} sebagai empat
digit heksadesimal, sehingga mahasiswa dapat membaca nilai segmen
dan offset yang sesungguhnya tanpa harus membuka emulator.

\subsection{Contoh 1: Menampilkan Offset Variabel}

Program pada Kode~\ref{lst:alamat-fisik} mendeklarasikan dua
variabel pada segmen data, lalu menampilkan offset keduanya dengan
operator \code{OFFSET}.

\lstinputlisting[
    language=MASM,
    caption={Program \file{alamat-fisik.asm} menampilkan offset
    variabel pada segmen data.},
    label={lst:alamat-fisik}
]{\CodeDir/alamat-fisik.asm}

Bagian yang perlu diperhatikan dari program tersebut adalah cara
nilai offset diperoleh:

\begin{lstlisting}[language=MASM, numbers=none]
MOV AX, OFFSET var_byte    ; AX = offset var_byte
MOV AX, OFFSET var_word    ; AX = offset var_word
\end{lstlisting}

Operator \code{OFFSET} menghasilkan jarak variabel tersebut dari
awal segmen data, bukan alamat fisiknya. Nilai inilah yang menjadi
bagian kedua pada pasangan \eng{segment:offset}. Karena variabel
\code{var\_byte} dideklarasikan lebih dahulu, offsetnya nol,
sedangkan \code{var\_word} berada tepat sesudahnya dengan offset
satu. Keluaran program tampak pada Kode~\ref{lst:keluaran-offset}.

\begin{session}[caption={Keluaran program \file{alamat-fisik.asm}.},
label={lst:keluaran-offset}]
Offset var_byte = 0000
Offset var_word = 0001
\end{session}

\begin{infobox}
Nilai offset yang tampil dapat berbeda satu atau dua byte bila
assembler yang dipakai menyelaraskan data ke alamat genap. Yang
penting bukan angkanya, melainkan caranya membaca: angka yang tampil
adalah offset, dan nomor segmennya harus dibaca terpisah dari
register \reg{DS}.
\end{infobox}

Untuk memperoleh alamat fisik, mahasiswa harus mengetahui nilai
\reg{DS} saat program berjalan. Misalkan pada Emu8086 atau DEBUG
terbaca \reg{DS} $= \hex{076A}$, maka kedua alamat fisik tersebut
dapat dihitung seperti pada Tabel~\ref{tab:analisis-offset}.

\begin{table}[H]
    \centering
    \caption{Perhitungan alamat fisik variabel pada program pertama
    dengan asumsi \reg{DS} $= \hex{076A}$.}
    \label{tab:analisis-offset}
    \small
    \begin{tabular}{@{}p{2.3cm}p{1.8cm}p{1.5cm}p{4.1cm}p{1.9cm}@{}}
        \toprule
        \textbf{Variabel} & \textbf{Offset} & \textbf{DS} &
        \textbf{Perhitungan} & \textbf{Alamat fisik} \\
        \midrule
        \code{var\_byte} & \hex{0000} & \hex{076A} &
        \hex{076A0} + \hex{0000} & \hex{076A0} \\
        \addlinespace
        \code{var\_word} & \hex{0001} & \hex{076A} &
        \hex{076A0} + \hex{0001} & \hex{076A1} \\
        \bottomrule
    \end{tabular}
\end{table}

\subsection{Contoh 2: Membaca Nilai Segmen dan Menghitung Alamat}

Program kedua, yang ditampilkan pada Kode~\ref{lst:hitung-alamat},
melangkah lebih jauh: ia membaca nomor segmen dengan operator
\code{seg}, membaca offset, lalu menghitung alamat fisiknya di dalam
program.

\lstinputlisting[
    language=MASM,
    caption={Program \file{hitung-alamat.asm} menghitung alamat
    fisik 20-bit dari pasangan segmen:offset.},
    label={lst:hitung-alamat}
]{\CodeDir/hitung-alamat.asm}

Inti perhitungannya terletak pada empat instruksi berikut.

\begin{lstlisting}[language=MASM, numbers=none]
MOV AX, seg var_nilai      ; AX = nomor segmen var_nilai
MOV DX, AX                 ; DX akan menerima bit 16..19
MOV CL, 4
SHR DX, CL                 ; DX = segmen / 16
SHL AX, CL                 ; AX = segmen x 16
ADD AX, BX                 ; AX = (segmen x 16) + offset
ADC DX, 0                  ; tambahkan carry ke DX
\end{lstlisting}

Perhatikan bahwa pergeseran ke kiri empat bit (\lstasm|SHL AX, CL|)
adalah bentuk lain dari perkalian dengan 16, dan pergeseran ke
kanan empat bit (\lstasm|SHR DX, CL|) memindahkan bagian segmen yang
melewati 16 bit ke register \reg{DX}. Karena hasil penjumlahan
dapat mencapai 20 bit, alamat fisik disimpan pada pasangan register
\reg{DX}:\reg{AX}, dengan \reg{DX} memuat bit 16 sampai 19 dan
\reg{AX} memuat bit 0 sampai 15. Keluaran program tampak pada
Kode~\ref{lst:keluaran-hitung}.

\begin{session}[caption={Keluaran program \file{hitung-alamat.asm}.},
label={lst:keluaran-hitung}]
Segmen: 076A
Offset: 0000
Alamat fisik: 000076A0
\end{session}

Keluaran tersebut menunjukkan tiga nilai sekaligus: nomor segmen
\hex{076A}, offset \hex{0000}, dan alamat fisik \hex{076A0} yang
ditulis delapan digit karena menempati pasangan \reg{DX}:\reg{AX}.
Mahasiswa dapat mengubah isi \reg{AX} sebelum perhitungan, memindahkan
variabel \code{var\_nilai} ke posisi lain pada segmen data, atau
mengganti segmen yang dipakai, lalu memeriksa apakah alamat fisik
yang dihasilkan tetap sesuai dengan perhitungan manual.

\begin{infobox}
Operator \code{seg} hanya menghasilkan \emph{nomor} segmen, bukan
alamat dasarnya. Pada contoh di atas, \lstasm|MOV AX, seg var_nilai|
memberi \hex{076A}, sedangkan alamat dasar segmen tersebut adalah
\hex{076A0}. Perbedaan satu digit nol inilah yang menjadi sumber
kesalahan paling sering pada perhitungan alamat.
\end{infobox}
